\subsection{Utilizzo della chat}

La chat \`e il cuore di SmartOrder: \`e lo strumento principale con cui puoi creare e gestire il tuo ordine.\\
Puoi scrivere in modo naturale, come se stessi parlando con un assistente, ad esempio chiedendo di aggiungere un prodotto, modificare una quantit\`a o rimuovere un articolo dal carrello.\\
Il sistema interpreta la tua richiesta e aggiorna automaticamente l'ordine.

\begin{figure}[H]
	\centering
	\includegraphics[width=0.8\textwidth]{ManualeUtente/chat}
	\caption{Area chat}
\end{figure}

\subsubsection{Inserimento messaggi}

Per inviare un messaggio, utilizza il campo di input situato nella parte inferiore della chat.

\begin{figure}[H]
	\centering
	\includegraphics[width=1\textwidth]{ManualeUtente/chat_barra}
	\caption{Barra di input della chat}
\end{figure}

Puoi scrivere messaggi fino a un massimo di 4096 caratteri.

Per inviare il messaggio hai due possibilit\`a:
\begin{itemize}[itemsep=5pt, parsep=5pt, label=$\scriptstyle\bullet$]
	\item Premere il tasto Invio sulla tastiera
	\item Cliccare sul pulsante di invio (1) accanto al campo di input
\end{itemize}

Se invece vuoi andare a capo senza inviare il messaggio, puoi usare la combinazione \texttt{Shift + Invio}.

Una volta inviato, il messaggio viene subito visualizzato nella chat e il sistema inizia ad elaborarlo.

\subsubsection{Funzionalit\`a avanzate}

Oltre ai messaggi di testo, puoi interagire con la chat anche utilizzando la voce.

\begin{itemize}[itemsep=5pt, parsep=5pt, label=$\scriptstyle\bullet$]
	\item \textbf{Invio di file audio}: puoi allegare un file audio cliccando sull'apposita icona (2) nel campo di input. Il sistema analizzer\`a il contenuto e lo convertir\`a automaticamente in una richiesta testuale.
\end{itemize}

I file supportati devono:
\begin{itemize}[itemsep=5pt, parsep=5pt, label=$\scriptstyle\bullet$]
	\item Avere una durata massima di 2 minuti
	\item Avere dimensione massima di 10 MB
	\item Essere nei formati \texttt{.mp3}, \texttt{.m4a}, \texttt{.flac}, \texttt{.wav}, \texttt{.webm}
\end{itemize}

Se il file non rispetta questi requisiti, ad esempio per formato non valido o dimensione eccessiva, il sistema mostrer\`a un messaggio di errore e il caricamento non verr\`a completato.

\begin{itemize}[itemsep=5pt, parsep=5pt, label=$\scriptstyle\bullet$]
	\item \textbf{Input vocale (microfono)}: in alternativa, puoi registrare direttamente un messaggio vocale cliccando sull'icona del microfono (3).
\end{itemize}

Il sistema trascrive automaticamente ci\`o che dici e lo utilizza come messaggio nella chat.\\
Anche in questo caso, la registrazione ha una durata massima di 2 minuti.

Se il limite viene superato, il sistema interrompe la registrazione e ti invita a riprovare con un messaggio pi\`u breve.

\subsubsection{Risposte del sistema}

Dopo aver inviato un messaggio, il sistema risponde automaticamente per aiutarti a costruire il tuo ordine.

Le risposte possono essere di diverso tipo:

\begin{itemize}[itemsep=5pt, parsep=5pt, label=$\scriptstyle\bullet$]
	\item \textbf{Conferme}: il sistema ti conferma le azioni eseguite, ad esempio quando un prodotto viene aggiunto o quando la quantit\`a viene aggiornata.
\end{itemize}

\begin{itemize}[itemsep=5pt, parsep=5pt, label=$\scriptstyle\bullet$]
	\item \textbf{Suggerimenti di prodotti}: se la richiesta \`e generica, ad esempio \vr{acqua}, il sistema ti propone una lista di prodotti tra cui scegliere.
\end{itemize}

\begin{itemize}[itemsep=5pt, parsep=5pt, label=$\scriptstyle\bullet$]
	\item \textbf{Richieste di chiarimento}: se il messaggio non \`e abbastanza preciso, il sistema ti chiede di specificare meglio cosa intendi, guidandoti nella scelta.
\end{itemize}

\begin{itemize}[itemsep=5pt, parsep=5pt, label=$\scriptstyle\bullet$]
	\item \textbf{Messaggi di errore}: se il sistema non riesce a interpretare la richiesta, ti verr\`a chiesto di riformularla in modo pi\`u chiaro.
\end{itemize}

Grazie a queste risposte, la chat ti accompagna passo dopo passo nella creazione dell'ordine.

\subsubsection{Stato del sistema}

Dopo aver inviato un messaggio, potresti notare un indicatore di caricamento nella chat.

Questo significa che il sistema sta elaborando la tua richiesta.\\
Una volta completata l'elaborazione, l'indicatore scompare e viene mostrata la risposta.
